\documentclass[12pt]{article}
\usepackage[margin=1in]{geometry}
\usepackage{amsmath,amsthm,amssymb,scrextend}
\usepackage{array}
\setlength{\parindent}{0pt}
\begin{document}
\newcommand{\ds}{\displaystyle}

\textbf{(i)}

$$P_2=\frac12$$

$$T_1\text{ sits in }S_1:\ T_3\text{ can sits in }S_3$$

$$T_1\text{ sits in }S_2:\ T_2\text{ chooses }S_1\text{ or }S_3$$

$$T_1\text{ sits in }S_3:\ T_3\text{ cannot sit in }S_3$$

$$P_3=\frac13+\frac13\cdot\frac12=\frac12$$

\bigskip\hrule\bigskip

\textbf{(ii)}

$$T_1\text{ sits in }S_1:\quad P(T_n\text{ sits in }S_n)=1$$
$$T_1\text{ sits in }S_n:\quad P(T_n\text{ sits in }S_n)=0$$

If $T_1$ sits in $S_k$, where $\ds 2\leq k\leq n-1$,


If $T_1$ sits in $S_k$, where $\ds 2\leq k\leq n-1$, $T_2,T_3,\ldots,T_{k-1}\text{ will sit in their own seats}$

Then $T_k$ is the first passenger whose seat is not available

There are $n-k+1$ available seats left: $\ds S_1,S_{k+1},S_{k+2},\ldots,S_n$
$$P(T_n\text{ sits in }S_n\mid T_1\text{ sits in }S_k)=P_{n-k+1}$$


This becomes the same problem with $n-k+1$ seats.

\begin{align*}
P_n
&=\frac1n\cdot 1+\frac1n\sum_{k=2}^{n-1}P_{n-k+1}+\frac1n\cdot 0\\
&=\frac1n\left(1+\sum_{k=2}^{n-1}P_{n-k+1}\right)\\
&=\frac1n\left(1+\sum_{r=2}^{n-1}P_r\right)
\end{align*}



\bigskip\hrule\bigskip

\textbf{(iii)}

Assume $\ds P_k=\frac12\quad (k\geq2)$

Base case: $\ds P_2=\frac12$


$$P_k=\frac1k\left(1+\sum_{r=2}^{k-1}P_r\right)\implies 1+\sum_{r=2}^{k-1}P_r=kP_k$$

\begin{align*}
P_{k+1}
&=\frac1{k+1}\left(1+\sum_{r=2}^{k}P_r\right)\\
&=\frac1{k+1}\left(1+\sum_{r=2}^{k-1}P_r+P_k\right)\\
&=\frac1{k+1}(kP_k+P_k)\\
&=\frac12
\end{align*}

By mathematical induction, $\ds P_n=\frac12$

\bigskip\hrule\bigskip

\textbf{(iv)}


$$Q_2=\frac12$$


For $n\geq3$, condition on the seat chosen by $T_1$.

$$T_1\text{ sits in }S_1 \text{ or }S_n:\ T_{n-1}\text{ sits in }S_{n-1}$$

$$T_1\text{ sits in }S_{n-1}:\ T_{n-1}\text{ cannot sit in }S_{n-1}$$

If $T_1$ sits in $S_k$, where $2\leq k\leq n-2$,

$$T_k\text{ becomes the first displaced passenger}$$

$$P(T_{n-1}\text{ sits in }S_{n-1}\mid T_1\text{ sits in }S_k)=Q_{n-k+1}$$

\begin{align*}
Q_n
&=\frac1n\left(2+\sum_{k=2}^{n-2}Q_{n-k+1}\right)\\
&=\frac1n\left(2+\sum_{r=3}^{n-1}Q_r\right)
\end{align*}

Base case: $\ds Q_3=\frac13+\frac13=\frac23$

Assume $\ds Q_3=Q_4=\cdots=Q_{n-1}=\frac23$

\begin{align*}
Q_n
&=\frac1n\left(2+\sum_{r=3}^{n-1}Q_r\right)\\
&=\frac1n\left(2+\frac{2(n-3)}{3}\right)\\
&=\frac23
\end{align*}

$$\therefore Q_n=
\begin{cases}
\ds \frac12, & n=2,\\[6pt]
\ds \frac23, & n\geq3.
\end{cases}$$

\bigskip
\hrule
\begin{center}
    This solution was submitted by \textbf{Jon}\\[4pt]
    Published and refined by \textbf{Jon}\\[8pt]
    \textit{Last updated: July 22, 2026}
\end{center}

\end{document}